% =====================================================================
%  1. OBJETIVO
% =====================================================================
\section{Objetivo}
Sustituir la funci\'{o}n \texttt{averageRangeInWindow()} de Lab-MR 1-II por una funci\'{o}n \texttt{extractWall()} que estime cada pared del pasillo como una recta $y = mx + b$ mediante \textbf{regresi\'{o}n lineal por m\'{i}nimos cuadrados} sobre los puntos l\'{a}ser. Crear el nodo \texttt{corr\_nav\_wall} (basado en el de la pr\'{a}ctica anterior), adaptar el bucle de control para usar las dos rectas en el c\'{a}lculo del punto objetivo de Pure Pursuit, compilar y lanzar con \texttt{ros2 launch seg\_tray corr\_nav\_wall\_launch.py}.


% =====================================================================
%  2. IMPLEMENTACIÓN
% =====================================================================
\section{Implementaci\'{o}n}

El nodo \texttt{corr\_nav\_wall\_node} se ha derivado del de Lab-MR 1-II reemplazando \texttt{averageRangeInWindow()} por \texttt{extractWall()} y eliminando el servicio que ya no se necesita. La estructura general (callback de l\'{a}ser, callback de odometr\'{i}a, bucle de control a 40~Hz que publica \texttt{cmd\_vel}) se mantiene. Las decisiones no obvias son:

\paragraph{Rangos angulares estrechos.} La pared izquierda se extrae del sector $[60^\circ, 120^\circ]$ y la derecha del sector $[-120^\circ, -60^\circ]$. Rangos m\'{a}s amplios (e.g. $\pm 30^\circ$ a $\pm 150^\circ$) incluyen lecturas de las esquinas del pasillo, que sesgan la regresi\'{o}n y hacen oscilar la orientaci\'{o}n estimada.

\paragraph{Filtrado de lecturas inv\'{a}lidas.} Se descartan rangos no finitos o no positivos antes de convertir a cartesianas. Si quedan menos de $2$ puntos v\'{a}lidos, la pared se marca como no detectada ($b = +\infty$).

\paragraph{Protecci\'{o}n frente a paredes verticales.} La f\'{o}rmula de m\'{i}nimos cuadrados para $m$ tiene como denominador $n\sum x_i^2 - (\sum x_i)^2$, que tiende a cero cuando los puntos comparten la misma $x$ (pared perpendicular al eje del robot). En ese caso se devuelve $m = 0$ y $b = \bar{y}$, lo que equivale a tratar la pared como horizontal y evita un \texttt{NaN}.

\paragraph{Punto objetivo anal\'{i}tico.} Conocidas $(m_L, b_L)$ y $(m_R, b_R)$, el centro del pasillo a la distancia de look-ahead $L$ es directamente
\begin{equation*}
  y_\text{goal} = \frac{m_L + m_R}{2}\, L + \frac{b_L + b_R}{2},
\end{equation*}
y la curvatura de Pure Pursuit es $\kappa = 2\,y_\text{goal} / L^2$. No hace falta promediar lecturas a izquierda y derecha como en Lab-MR 1-II: el centrado y la orientaci\'{o}n quedan absorbidos por el ajuste lineal.

\paragraph{Cl\'{a}usula de guarda al perder una pared.} Si $b_L$ o $b_R$ no son finitos (e.g. en una curva el sector angular deja de ver una pared), se publica $v = \omega = 0$ en vez de calcular comandos err\'{a}ticos. Es una protecci\'{o}n de seguridad m\'{a}s que un modo de curva.


% =====================================================================
%  3. PRUEBAS
% =====================================================================
\section{Pruebas}

Tras compilar e incluir el ejecutable en \texttt{CMakeLists.txt} y crear el launch file copiado del de Lab-MR 1-II, el nodo se lanza y el robot navega por el centro del pasillo en el escenario de \coppelia{} (Figura~\ref{fig:coppelia}).

\begin{figure}[H]
  \centering
  \includegraphics[width=0.85\linewidth]{coppelia_resultado.png}
  \caption{Resultado de la navegaci\'{o}n con estimaci\'{o}n de paredes en \coppelia{}: el robot \pioneer{} sigue el pasillo de forma centrada usando los modelos de pared izquierda y derecha.}
  \label{fig:coppelia}
\end{figure}

En tramos rectos, el comportamiento es notablemente m\'{a}s suave que con el m\'{e}todo de ventanas angulares de Lab-MR 1-II porque la regresi\'{o}n integra todas las lecturas del sector y filtra el ruido individual. La orientaci\'{o}n del robot se obtiene directamente de las pendientes $m_L$ y $m_R$, sin necesidad de la heur\'{i}stica de comparar lecturas a $\pm 15^\circ$ que us\'{a}bamos antes. La velocidad de crucero se ha podido subir a $\SI{0.5}{\meter\per\second}$ frente a los $\SI{0.25}{\meter\per\second}$ de la pr\'{a}ctica anterior, gracias a la mayor estabilidad del estimador.

En las curvas, esta versi\'{o}n no implementa un modo dedicado: cuando una pared sale del sector de detecci\'{o}n el intercepto se vuelve infinito y la cl\'{a}usula de guarda detiene el robot. La integraci\'{o}n con el modo curva de Lab-MR 1-II queda como mejora natural pero no la pide el enunciado de esta pr\'{a}ctica.


% =====================================================================
%  4. COMENTARIOS FINALES
% =====================================================================
\section{Comentarios finales}

La sustituci\'{o}n de \texttt{averageRangeInWindow()} por \texttt{extractWall()} cumple lo que pide el enunciado: las paredes se modelan como rectas $y = mx + b$ por m\'{i}nimos cuadrados, el punto objetivo se calcula anal\'{i}ticamente como el centro del pasillo a la distancia de look-ahead, y el control mantiene Pure Pursuit. La estimaci\'{o}n por regresi\'{o}n es m\'{a}s robusta al ruido que el promedio en ventanas estrechas, lo que se traduce en trayectorias m\'{a}s suaves y velocidad de crucero m\'{a}s alta en tramos rectos.
